\documentclass[11pt]{article}
\usepackage[margin=0.75in]{geometry}
\usepackage{amsmath,amssymb,amsthm}
\usepackage{array,booktabs,xcolor,colortbl}
\usepackage{tikz}
\usepackage{hyperref}
\setlength{\parindent}{0pt}
\newcommand{\stepnum}[1]{%
  \tikz[baseline=(n.base)] \node[draw,circle,inner sep=1.6pt,
    font=\scriptsize\bfseries,fill=white] (n) {#1};}%
\newtheorem{theorem}{Theorem}
\theoremstyle{definition}
\newtheorem{definition}{Definition}
\title{Flow Proof Artifact: matmul\_vectorize}
\author{Generated by \texttt{flow doc proof}}
\date{\today}
\begin{document}
\maketitle
\bigskip
\noindent{\large \textbf{Derived fact 10.} \textit{Matmul\_vectorized\_correct} \hfill \texttt{matmul\_vectorized\_correct}}\par
\medskip
\noindent\fcolorbox{black!8}{%
\begin{minipage}{\dimexpr\textwidth-2\fboxsep-2\fboxrule}
\textbf{Goal.} We're showing that matrices\_equal(C\_naive, C\_fast, m, n).\\[0.35em]
$\displaystyle \forall m \in \mathbb{Z},\; matrices_equal(C_naive, C_fast, m, n)$
\end{minipage}%
}\par
\medskip
\noindent
\begin{tabular}{@{} c >{\raggedright\arraybackslash}p{0.40\textwidth} @{\hspace{0.4em}\color{gray!50}\vrule width 0.5pt\hspace{0.4em}} c >{\raggedleft\arraybackslash}p{0.40\textwidth} @{}}
\multicolumn{2}{l}{\textbf{Proof}} & \multicolumn{2}{r}{\textbf{Mathematics}} \\
\midrule
\stepnum{1} & Let A = arbitrary\_f32\_matrix(m, k). & & \multicolumn{1}{p{0.40\textwidth}}{} \\[0.55em]
\stepnum{2} & Let B = arbitrary\_f32\_matrix(k, n). & & \multicolumn{1}{p{0.40\textwidth}}{} \\[0.55em]
\stepnum{3} & Let C\_naive = f32\_matrix\_zeros(m, n). & & \multicolumn{1}{p{0.40\textwidth}}{} \\[0.55em]
\stepnum{4} & Let C\_fast = f32\_matrix\_zeros(m, n). & & \multicolumn{1}{p{0.40\textwidth}}{} \\[0.55em]
\stepnum{5} & From ①, ②, ③, and ④, this implies matrices\_equal(C\_naive, C\_fast, m, n). Hence proven. & \stepnum{5} & $matrices_equal(C_naive, C_fast, m, n)$ \\[0.55em]
\bottomrule
\end{tabular}
\label{thm:matmul_vectorized_correct}
\par\medskip\hrule
\bigskip
\noindent{\large \textbf{Derived fact 11.} \textit{Loop\_fusion\_correct} \hfill \texttt{loop\_fusion\_correct}}\par
\medskip
\noindent\fcolorbox{black!8}{%
\begin{minipage}{\dimexpr\textwidth-2\fboxsep-2\fboxrule}
\textbf{Goal.} We're showing that memory\_equal(σ\_separate, σ\_fused).\\[0.35em]
$\displaystyle \forall n \in \mathbb{Z},\; memory_equal(σ_separate, σ_fused)$
\end{minipage}%
}\par
\medskip
\noindent
\begin{tabular}{@{} c >{\raggedright\arraybackslash}p{0.40\textwidth} @{\hspace{0.4em}\color{gray!50}\vrule width 0.5pt\hspace{0.4em}} c >{\raggedleft\arraybackslash}p{0.40\textwidth} @{}}
\multicolumn{2}{l}{\textbf{Proof}} & \multicolumn{2}{r}{\textbf{Mathematics}} \\
\midrule
\stepnum{1} & Let σ = arbitrary\_memory(). & & \multicolumn{1}{p{0.40\textwidth}}{} \\[0.55em]
\stepnum{2} & Let a = fresh\_array(n). & & \multicolumn{1}{p{0.40\textwidth}}{} \\[0.55em]
\stepnum{3} & Let b = fresh\_array(n). & & \multicolumn{1}{p{0.40\textwidth}}{} \\[0.55em]
\stepnum{4} & Let σ\_separate = run\_separate\_loops(σ, a, b, n). & & \multicolumn{1}{p{0.40\textwidth}}{} \\[0.55em]
\stepnum{5} & Let σ\_fused = run\_fused\_loop(σ, a, b, n). & & \multicolumn{1}{p{0.40\textwidth}}{} \\[0.55em]
\stepnum{6} & From ①, ②, ③, ④, and ⑤, this implies memory\_equal(σ\_separate, σ\_fused). Hence proven. & \stepnum{6} & $memory_equal(σ_separate, σ_fused)$ \\[0.55em]
\bottomrule
\end{tabular}
\label{thm:loop_fusion_correct}
\par\medskip\hrule
\end{document}
